We study a central question in LLM reliability: when commonsense reasoning fails, is the main issue missing knowledge or a breakdown in integrating question goals with candidate actions? We present a paired behavioral--mechanistic protocol that evaluates the same \csqa examples under \direct and \cotprompt prompting, then applies causal-style interventions in a local open-model proxy. On 120 validation items, \gptapi reaches 90.0\% accuracy with \direct prompting and 89.2\% with \cotprompt prompting, with low bidirectional failure flips (\direct$\rightarrow$\cotprompt: 2.5\%; \cotprompt$\rightarrow$\direct: 1.7\%). A paired McNemar exact test shows no significant difference ($p=1.0$). We further run layer-wise residual ablations in \proxy on matched success/failure examples and test whether a hypothesized target layer has stronger causal effect than control layers. The targeted effect is not statistically significant (paired $t$-test $p=0.853$, Wilcoxon $p=0.546$, Cohen's $d=0.038$). These results do not support a strong claim that \cotprompt broadly degrades commonsense performance in this sampled setting, and they provide only weak evidence for localized coherence circuits with this small proxy model. The main practical takeaway is methodological: robust parsing and paired evaluation design are necessary before drawing mechanistic conclusions from apparent reasoning failures.
